\section{Conclusão e Trabalhos Futuros}
\label{sec:conclusao}

Este trabalho avaliou como a preservação do estado e a frequência de falhas afetam a recuperação de sessões e os efeitos das leituras em redes IoT, comparando configurações clássica, híbrida e pós-quântica. No protótipo, os modos com chaves novas não apresentaram reutilizações de chave e nonce nos cenários testados. Nos experimentos TLS com deduplicação ativa, reinícios e confirmações perdidas em todas as leituras elegíveis e perda de pacotes configurada em 0\%, a configuração pós-quântica apresentou economia de 67,87\% nos bytes IPv4 por leitura confirmada com contexto preservado e aumento de 14,24\% com contexto perdido, frente à autenticação completa. Esses percentuais correspondem às cadeias avaliadas, que incluem a CA raiz. A dependência da política de estado também foi observada nas configurações clássica e híbrida. Na matriz principal, a deduplicação resultou em zero efeitos duplicados no banco, contra 2.700 sem supressão. Os resultados mostram que o benefício da retomada depende do contexto disponível, enquanto o controle de duplicações exige preservar o histórico da aplicação.

Como trabalhos futuros, propõe-se avaliar cadeias de certificados sem transmissão da CA raiz e comparar tickets que transportam estado com referências a sessões armazenadas no gateway, mantendo a preservação do contexto como variável separada. A avaliação em sensores físicos permitirá medir memória, tempo de recuperação e consumo energético, enquanto múltiplos sensores e diferentes cargas permitirão investigar a concorrência. Também se pretende implementar uma fila persistente de leituras pendentes e investigar a proteção do histórico do gateway contra restauração de versões antigas. Por fim, a extensão para atuadores exige coordenar o registro persistente da decisão com a execução do efeito externo.
